\section{Aggregation and Dissolution}
\label{sec:pathways}

\begin{definition}[Critical contact number]
\derived\ $N_c^*=3$. A regular tetrahedron requires four centres; three contacts
are necessary and sufficient for one particle to complete a new tetrahedral unit.
The same count attaches a new tetrahedron to an exposed face of a Boerdijk--Coxeter
helix.
\end{definition}

\begin{itemize}
\item If $N_c<3$ the particle follows the \textbf{dissolution pathway}: ordered
layering is lost and capacity returns to the ambient field.
\item If $N_c\ge 3$ the particle follows the \textbf{aggregation pathway}: residual
tension couples contacting membranes into larger composite structure.
\end{itemize}

Nucleation of a new cell is geometrically natural at a charged opposing ($P$ against
$p$) interface when residual accumulation near critical shortfall discharges into
the solvent gap. Dissolution is open for under-coordinated particles ($N_c<3$),
with elevated rate on the periphery and under solitary caustic spikes.
